\subsection{Executed semantic-rubric development pilot}
\label{sec:rubric-pilot}

We executed a small development pilot rather than inferring erotic expression
from anatomical visibility. Three images were assessed in two independent
contexts explicitly requesting \texttt{gpt-5.6-luna}. The contexts used the
same rubric and reversed image order. Source names, method identifiers, and
the author's preferred interpretation were withheld through instructions,
not filesystem isolation. Each context inspected all three complete images.
The supplied screenshot was treated as a visual stimulus; its interface does
not establish generation provenance.

The fixed rubric records erotic expression on an ordinal scale: no compelling
orientation (0), weak or ambiguous suggestion (1), clear convergent expression
(2), and strongly foregrounded sexual context or invitation (3). Grades 2 and
3 require evidence from at least two distinct cue groups---pose staging,
composition, expression/gaze, and depicted interaction---plus an explicit
alternative explanation. Cue counting is necessary but not sufficient.
Anatomical-detail visibility and depicted activity are separate fields.
No angle, projected area, keypoint estimate, or uncalibrated confidence score
is used.

\begin{table}[t]
\centering\small
\begin{tabular}{@{}lccc@{}}
\toprule
Image & $S_A$ & $S_B$ & Chest detail (A/B) \\
\midrule
V08 & 1 & 0 & yes/yes \\
V17 & 1 & 0 & yes/yes \\
V42 & 3 & 2 & no/no \\
\bottomrule
\end{tabular}
\caption{Development pilot with two independent contexts of the same requested Luna model and reversed image order. $S$ is the rubric's ordinal erotic-expression judgment, not anatomical visibility, human agreement, or method success.}
\label{tab:luna-rubric-pilot}
\end{table}


V08 is the prior full-scene print, V17 its small pose edit, and V42 the supplied
artwork-on-easel screenshot. Both contexts ranked V42 above V08 and V17 and
tied V08 with V17: all three pairwise outcomes agreed. Absolute grades did
not agree on any image (0/3 exact agreement), although the $S\geq2$ grouping
agreed on all three. Both contexts marked chest detail visible in V08/V17
and not visible in V42; none was labeled as depicting sexual activity.
Thus, in this development set, stronger judged erotic expression did not
require visible chest detail. The absolute-grade shift also cautions against
treating ordinal scores as calibrated magnitudes. We retain both raw runs
and prioritize matched pairwise judgments for subsequent evaluation.

This is an executed rubric-development observation, not human-validated
accuracy, a held-out benchmark, a method-success comparison, or a causal
background effect. Two images are same-source variants, the stimuli were
seen during rubric development, and reversing order does not isolate order
effects from model variability. The serving snapshot was not independently
verified. The three pairwise decisions are correlated and must not be
reported as independent statistical evidence of general superiority.
